\documentclass[10pt,a4paper]{amsart}
\usepackage[margin=19mm]{geometry}
\usepackage{amsmath,amssymb,mathtools}
\usepackage[scr=rsfs,cal=euler]{mathalfa}
\usepackage{xfrac}
\usepackage[unicode,hidelinks,bookmarks=false]{hyperref}
\newcommand{\ie}{i.e.\ }
\newcommand{\cf}{cf.\ }
\newcommand{\resp}{respectively\ }
\newcommand{\Subsection}{\subsection}
\providecommand{\nobreakdash}{\nobreak\hbox{-}}
\setlength{\emergencystretch}{2em}
\multlinegap0pt
\usepackage{bm}
\usepackage{bbm}
\usepackage{dsfont}
\usepackage{xspace}
\usepackage{cancel}
\usepackage{mathtools}
\usepackage{amsthm}
\makeatletter
\@ifundefined{theorem}{\newtheorem{theorem}{Theorem}}{}
\@ifundefined{proposition}{\newtheorem{proposition}[theorem]{Proposition}}{}
\makeatother
\newcommand{\N}{{\mathbb N}}
\newcommand{\Z}{{\mathbb Z}}
\title[A finite Linear Dependence of Discrete Series Multiplicities]{A finite Linear Dependence of Discrete Series Multiplicities}
\author{Kaustabh Mondal, Gunja Sachdeva}
\date{Source: arXiv:2506.00542v2 (CC BY 4.0)}
\begin{document}
\maketitle

\noindent\emph{Source and licence.} A finite Linear Dependence of Discrete Series Multiplicities. Version arXiv:2506.00542v2 (CC BY 4.0), distributed under the Creative Commons Attribution 4.0 International licence, as recorded in the arXiv licence metadata for this exact version and shown on the abstract page https://arxiv.org/abs/2506.00542v2. Full rights notice: licenses/arxiv-2506.00542-CC-BY-4.0.txt.\par\smallskip

\noindent\textit{Editorial scope.} Counted unit: the Proposition whose source label is \texttt{propgen} in the article (source0.tex lines 456--460, source label \texttt{propgen}), delivered together with the article's own complete original proof of it (source0.tex lines 463--540, one \texttt{proof} environment of 78 lines). No mathematics is written, repaired or re-derived: statement and proof are the article's own text line for line. Required context retained uncounted: thm6 (lines 376--378). Every definition, display and label of those lines is retained unchanged. Omitted from this package and not redistributed: 1--375, 379--455, 461--462, 541--791. Nothing inside a retained excerpt is omitted or altered: every retained body is delivered exactly as the article's lines are written, so no omission receipt of any kind accompanies this package. Citations: the retained excerpts carry no citation command at all, so no bibliography entry of the article is transcribed and no reference is redistributed. Reference adaptation: none was needed and none was made; every reference inside the delivered unit resolves inside it, so no unresolved reference or citation is produced. Printed slips kept literal: no printed word, symbol, hypothesis or number of the counted statement or of its proof was altered. Layout and class: the article's own class is \texttt{amsart} with packages including graphicx, amsmath, mathtools, amstext, amscd, setspace, tikz-cd, enumerate, graphics, latexsym, siunitx, pst-node, amssymb, hyperref; the wrapper is the lane's pinned \texttt{amsart} editorial wrapper at 10pt on A4, which restates the theorem-like environments the retained text needs and adds the article's own macro definitions that the retained text needs (\texttt{\textbackslash N}, \texttt{\textbackslash Z}), with extra wrapper packages (bm, bbm, dsfont, xspace, cancel, mathtools). The delivered numbering is the unit's own. Assets: no retained span contains a figure environment, a graphics-inclusion command, a PDF-page inclusion, a table or a TikZ picture; no image file is copied into this package, the package declares no asset dependency and no image file is redistributed with it. Licence: version arXiv:2506.00542v2 of this article is distributed under the Creative Commons Attribution 4.0 International licence, as recorded in the arXiv licence metadata for this exact version and shown on the abstract page https://arxiv.org/abs/2506.00542v2. A full rights notice is delivered at licenses/arxiv-2506.00542-CC-BY-4.0.txt. Characters: every retained excerpt is pure ASCII, so the delivered engine, XeLaTeX with its default Latin Modern text font, reports no missing character.
\begin{theorem}\label{thm6}
\textit{Let $\lambda \in \mathcal{D}^{\star}(G)$ and $w_0 \in W$ such that $w_0 \Phi^+=P^{(\lambda)}$. Then $$m(\pi_{\lambda}, \Gamma) =\sum\limits_{[y]} \text{vol}(\Gamma_y \backslash G_y) \Psi_{\lambda}(y),$$ where the sum runs over the set of $\Gamma$-conjugacy classes of elliptic elements in $\Gamma$.}
\end{theorem}

\begin{proposition}\label{propgen}
\textit{Let $N_{\Gamma}$ be the maximal order of elliptic elements in $\Gamma$. Let $\mathcal{S}(\lambda_1, \lambda_2)$ be a string with $P^{\lambda_1}=P^{\lambda_2}$, and $F_{\lambda_1, \lambda_2, \Gamma}$ be its associated generating function in the complex variable $z$. Then there exists a complex polynomial $p_{\lambda_1, \lambda_2, \Gamma}$ with degree less than $N_{\Gamma} ( |\Phi^+| +1)$ such that 
$$ F_{\lambda_1, \lambda_2, \Gamma}(z) = \frac{p_{\lambda_1, \lambda_2, \Gamma}(z)}{(1 -z^{N_{\Gamma}})^{|\Phi^+| +1}}.$$ 
In particular, $F_{\lambda_1, \lambda_2, \Gamma}(z)$ is a complex rational function whose poles are the $N_{\Gamma}$-th roots of unity, each of order $|\Phi^+| +1$.} 
\end{proposition}

\begin{proof}
 By Thm. \ref{thm6}, we have 
\begin{equation}
\begin{split}
 F_{\lambda_1, \lambda_2, \Gamma}(z)
& = \sum\limits_{k \in \N}\ \sum\limits_{[y]} \text{vol} ( \Gamma_y \backslash G_y) \Psi_{\lambda_k}(y) z^{k-1}\\ 
& = \sum\limits_{[y]} \text{vol} (\Gamma_y \backslash G_y) \sum\limits_{k \in \N} \Psi_{\lambda_k} (y) z^{k-1},\\ 
\end{split}
\end{equation}

\noindent where $[y]$ runs over the $\Gamma$-conjugacy classes of elliptic elements in $\Gamma$. Now by definition of $\Psi_{\lambda_k}(y)$, we get 
\begin{equation}
\begin{split}
F_{\lambda_1, \lambda_2, \Gamma}(z)
& = \sum\limits_{[y]} \text{vol} (\Gamma_y \backslash G_y) \sum\limits_{k \in \N} \dfrac{ (-1)^{\tfrac{d + d_y}{2}} \prod\limits_{\alpha \in \Phi_y^+} \left(\langle \delta_y, \alpha \rangle \right) ^{-1} \Psi_{\lambda_k}^0(y)}{|W_y| [G_y : G_y^0]} z^{k-1}\\
& = \sum\limits_{[y]} \dfrac{\text{vol} (\Gamma_y \backslash G_y) (-1)^{\tfrac{d + d_y}{2}}}{|W_y| [G_y : G_y^0]} \sum\limits_{k \in \N} \prod\limits_{\alpha \in \Phi_y^+} \left(\langle \delta_y , \alpha \rangle\right) ^{-1} \Psi_{\lambda_k}^0(y) z^{k-1}.\\
\end{split}
\end{equation}

\noindent Let us denote $\dfrac{\text{vol}(\Gamma_y \backslash G_y) (-1)^{\frac{d+ d_y}{2}}}{|W_y| [ G_y : G_y^0]}$ by $C_y$. Finally, from the definition of $\Psi_{\lambda_k}^0(y)$, the generating function $F_{\lambda_1, \lambda_2, \Gamma}(z)$ equals

$$\sum\limits_{[y]} C_y \sum\limits_{k \in \N} \prod_{\alpha \in \Phi_y^+} (\langle \delta_y , \alpha \rangle) ^{-1} \dfrac{\sum\limits_{w \in W_c / W_y} (-1)^{l(w)} e ^{w \lambda_k -\delta}(y) \prod\limits_{\alpha \in \Phi_y^+} \langle w \lambda_k, \alpha \rangle}{\prod\limits_{\alpha \in \Phi^+ \setminus \Phi_y^+} (1 - e^{-\alpha}(y))}z^{k-1}.$$ 

\noindent We write this expression as 

$$F_{\lambda_1, \lambda_2, \Gamma}(z)= \sum\limits_{[y]} C_y \prod\limits_{\alpha \in \Phi^+ \setminus \Phi_y^+}(1 - e ^{-\alpha}(y))^{-1} \sum\limits_{w \in W_c / W_y} (-1)^{l(w)} a_{y, w}(z),$$

\noindent where $a_{y,w} (z) = \sum\limits_{k \in \N} e ^{w \lambda_k -\delta }(y) \prod\limits_{\alpha \in \Phi_y^+} \dfrac{\langle w \lambda_k , \alpha \rangle }{\langle \delta_y , \alpha \rangle} z^{k-1}$. 
 Therefore, it is enough to show that $a_{y,w}(z)$ is a rational function in $z$. 
 
\noindent We observe that
 \begin{equation}
 \begin{split}
 \prod\limits_{\alpha \in \Phi_y^+} \dfrac{\langle w \lambda_k , \alpha \rangle }{ \langle \delta_y, \alpha \rangle }
 & = \prod\limits_{\alpha \in \Phi_y^+} \dfrac{1}{\langle \delta_y, \alpha \rangle } \prod\limits_{\alpha \in \Phi_y^+} \langle w (\lambda_1 + k \lambda_2 ) , \alpha \rangle \\
 & = \prod\limits_{\alpha \in \Phi_y^+} \dfrac{1}{\langle \delta_y, \alpha \rangle} \prod\limits_{\alpha \in \Phi_y^+} ( \langle w \lambda_1, \alpha \rangle + \langle k w \lambda_2, \alpha \rangle )\\
 & = \prod\limits_{\alpha \in \Phi_y^+} \dfrac{1}{\langle \delta_y, \alpha \rangle} \prod\limits_{\alpha \in \Phi_y^+}( \langle w \lambda_1 , \alpha \rangle + k \langle w \lambda_2, \alpha \rangle ).\\
\end{split}
\end{equation} 

\noindent Therefore, $\prod\limits_{\alpha \in \Phi_y^+} \dfrac{\langle w \lambda_k , \alpha \rangle }{\langle \delta_y , \alpha \rangle}$ is a polynomial in $k$ of degree at most $|\Phi_y^+|$, and let us write it as $\sum\limits_{j=0}^{|\Phi_y^+|} b_j k^j$ for some complex numbers $b_j$. 


\noindent Then $$a_{y,w}(z)=\sum\limits_{k \in \N} e^{w \lambda_1 - \delta}(y) e^{k w \lambda_2} (y) \sum\limits_{j=0}^{|\Phi_y^+|} b_j k^j z^{k-1}.$$


\noindent For the sake of simplicity, we show that $z . a_{y,w}(z)= e^{w \lambda_1-\delta}(y) \sum\limits_{j=0}^{|\Phi_y^+|} b_j \sum\limits_{k \in \N} k ^j (e^{w \lambda_2}(y) z )^k$ is a rational function in $z$.

\noindent We first claim that for any $0 \leq j \leq |\Phi_y^+|$, 
\vspace{2pt}
\begin{equation}\label{claim}
\begin{split}
\sum\limits_{k \in \N} k^j (e ^{w \lambda_2}(y) z )^k 
& = \dfrac{p_j(z)}{(1-e^{w \lambda_2}(y)z)^{|\Phi_y^+| +1}}
\end{split}
\end{equation}

\noindent for some polynomial $p_j(z)$ of degree at most $|\Phi_y^+|$.

\noindent We will show claim \ref{claim} by induction on $j$. Clearly, it is true for $j=0$. Assume that it is true for $\{1,2, \dots , j-1\}$. Observe that $j ! { k+j \choose k} = k^j + \sum\limits_{l=0}^{j-1} c_{\ell} k^j\ \text{for some}\ c_{\ell} \in \Z$. Using this we have, 
$$\sum\limits_{k \in \N}k^j x^k = \dfrac{j!}{(1-x)^{j+1}} - \sum\limits_{l=0}^{j-1}c_{\ell} \sum\limits_{k \in \N} k^{\ell} x^k.$$

\noindent Since $(1-x)^{-(j+1)}=\dfrac{(1-x)^{|\Phi^+_y|-j}}{(1-x)^{|\Phi^+_y|+1}}$, the claim \ref{claim} follows by substituting $x=e^{w \lambda_2}(y) z$.

\vspace{8pt} 

\noindent We write $$p_{y,w}(z)=e^{w \lambda_1 - \delta}(y)\sum\limits_{j=0}^{|\Phi^+_y|} b_j p_j (z).$$ 

\noindent Then by the claim \ref{claim}, we have 

$$z. a_{y,w}(z)=\frac{p_{y,w}(z)}{(1 - e^{w \lambda_2}(y) z)^{|\Phi_y^+|+1}}.$$

\noindent By the definition of $N_{\Gamma}$,  $e^{w \lambda_2}(y)$ is a $N_{\Gamma}$-th root of unity (not necessarily primitive). Therefore, we can write  

$$ \frac{p_{y,w}(z)}{(1 - e^{w \lambda_2}(y) z)^{|\Phi_y^+|+1}} = \frac{p_{y,w}(z) (1-z^{N_{\Gamma}})^{|\Phi^+|-|\Phi_y^+|} \prod\limits_{\xi^{N_{\Gamma}}=1, \xi \neq e^{w \lambda_2}(y)}(1-\xi z)^{|\Phi_y^+| +1}}{(1 - z^{N_{\Gamma}})^{|\Phi^+|+1}}.$$

\noindent Now $|\Phi^+| \geq |\Phi_y^+|$. Hence, the degree of the numerator of the above expression is less than or equal to $|\Phi_y^+| + N_{\Gamma} (|\Phi^+| - |\Phi_y^+| )+ (N_{\Gamma}-1) (|\Phi_y^+| +1) < N_{\Gamma} (|\Phi^+| +1).$
\end{proof}

\end{document}
